\documentclass[12pt]{article}

\input{../../_assets/preambulo.tex}


\begin{document}

    % 1. Foto de fondo
    % 2. Título
    % 3. Encabezado Izquierdo
    % 4. Color de fondo
    % 5. Coord x del titulo
    % 6. Coord y del titulo
    % 7. Fecha

    
    \input{../../_assets/portada}
    \portadaExamen{ffccA4.jpg}{Topología I\\Examen XV}{Topología I. Examen XV}{MidnightBlue}{-8}{28}{2026}{Roxana Acedo Parra}

    \begin{description}
        \item[Asignatura] Topología I.
        \item[Curso Académico] 2025/2026.
        \item[Grado] Doble Grado en Ingeniería Informática y Matemáticas.
        \item[Grupo] Único.
        \item[Profesor] Antonio Alarcón López.
        \item[Descripción] Convocatoria Extraordinaria.
        \item[Fecha] 9 de febrero de 2026.
        \item[Duración] 180 minutos.
    
    \end{description}
    \newpage


    % ------------------------------------
    
    \begin{ejercicio}
        En $X = \left[0, 1\right] \cup \left[2, 3\right] \subset \mathbb{R}$, se considera la familia de subconjuntos
        
        $$\mathcal{B} = \{\{0\}\} \cup \left\{ \left]x - \eps, x + \eps\right[ \;\middle|\; 0 < x < 1,\; 0 < \eps \leq \min\{x, 1 - x\} \right\} \cup $$ 
        $$\left\{ \left]1 - \eps, 1\right] \;\middle|\; 0 < \eps \leq 1 \right\} \cup \left\{ \{0, x\} \;\middle|\; x \in \left[2, 3\right] \right\}.$$
        
        \begin{enumerate}[label=\alph*)]
            \item Demuestra que existe una única topología $\mathcal{T}$ en $X$ tal que $\mathcal{B}$ es base de $(X, \mathcal{T})$.
            \item Calcula la adherencia y el interior en $(X, \mathcal{T})$ de los siguientes subconjuntos $B_1 = \{0\}$ y $B_2 = X \cap \mathbb{Q}$.
            \item ¿Es $(X, \mathcal{T})$ $T_1$? ¿Es $T_2$? ¿Es $1\text{AN}$? ¿Es $2\text{AN}$?
            \item ¿Es $(X, \mathcal{T})$ conexo? ¿Es compacto?
            \item Se considera el conjunto $A = \{0\} \cup \left[2, 3\right]$. Determina si el espacio topológico cociente $(X/A, \mathcal{T}/A)$ es homeomorfo a $\left(\left[0, 1\right], \mathcal{T}_u|_{\left[0, 1\right]}\right)$, donde $\mathcal{T}_u$ es la topología usual en $\mathbb{R}$.
            \item Se considera el conjunto $E = \{0, 1\} \cup \left[2, 3\right]$. Determina si el espacio topológico cociente $(X/E, \mathcal{T}/E)$ es homeomorfo a $\left(\left]0, 1\right], \mathcal{T}_u|_{\left]0, 1\right]}\right)$, donde $\mathcal{T}_u$ es la topología usual en $\mathbb{R}$.
        \end{enumerate}
    \end{ejercicio}

    \begin{ejercicio}
        Un espacio topológico $(X, \mathcal{T})$ se dice $T_3$ si es $T_1$ y cumple que para cada subconjunto cerrado $C \in \mathcal{C}_{\mathcal{T}}$ y cada punto $x \in X \setminus C$ existen abiertos $A_1, A_2 \in \mathcal{T}$ tales que $C \subset A_1$, $x \in A_2$ y $A_1 \cap A_2 = \emptyset$.
        \begin{enumerate}[label=\alph*)]
            \item Sea $(X, \mathcal{T})$ un espacio topológico $T_1$. Demuestra que $(X, \mathcal{T})$ es $T_3$ si, y solo si, para cada punto $x \in X$ y cada abierto $U \in \mathcal{T}$ con $x \in U$, existe un abierto $V \in \mathcal{T}$ tal que $x \in V$ y $\overline{V} \subset U$.
            \item Demuestra que $T_3$ es un invariante topológico.
        \end{enumerate}
    \end{ejercicio}

    \newpage
    \setcounter{ejercicio}{0} % Reiniciar contador de ejercicios
    \noindent

        \begin{ejercicio}
        En $X = \left[0, 1\right] \cup \left[2, 3\right] \subset \mathbb{R}$, se considera la familia de subconjuntos
        
        $$\mathcal{B} = \{\{0\}\} \cup \left\{ \left]x - \eps, x + \eps\right[ \;\middle|\; 0 < x < 1,\; 0 < \eps \leq \min\{x, 1 - x\} \right\} \cup $$ 
        $$\left\{ \left]1 - \eps, 1\right] \;\middle|\; 0 < \eps \leq 1 \right\} \cup \left\{ \{0, x\} \;\middle|\; x \in \left[2, 3\right] \right\}.$$
        
        \begin{enumerate}[label=\alph*)]
            \item Demuestra que existe una única topología $\mathcal{T}$ en $X$ tal que $\mathcal{B}$ es base de $(X, \mathcal{T})$.

                Veamos que la familia $\mathcal{B}$ verifica
    
                \textbf{B1:} $X$ puede escribirse como unión de elementos de $\mathcal{B}$. Esto es sencillo de comprobar ya que
                \[
                X=(0,1]\cup\left(\bigcup_{x\in[2,3]}\{0,x\}\right)
                \]
        
                \textbf{B2:} Sean $B, B'\in\mathcal{B}$ y $x\in B\cap B'$ existe $\tilde{B}\in\mathcal{B}$ tal que $x\in\tilde{B}\subseteq B\cap B'$.
    
                Veamos qué elementos de $\mathcal{B}$ pueden tener intersección.
                \begin{itemize}[leftmargin=*]
                    \item El elemento $\{0\}$ solo tiene intersección con los elementos de la forma $\{0,x\}$ con $x\in[2,3]$. En este caso $\{0\}\cap\{0,x\}=\{0\}\in\mathcal{B}$.
                    
                    Si $B=\{0,x\}$ y $B'=\{0,y\}$ con $x,y\in[2,3]$ tenemos $\{0,x\}\cap\{0,y\}=\{0\}\in\mathcal{B}$.
                    
                    \item Supongamos ahora que $B=(1-\lambda,1]$ y $B'=(1-\lambda',1]$ para $0<\lambda\le\lambda'\le 1$. Entonces se tiene
                    \[
                    (1-\lambda,1]\cap(1-\lambda',1]=(1-\lambda,1]\in\mathcal{B}
                    \]
        
                    \item Supongamos que $B=(x-\epsilon,x+\epsilon)$ para $0<x<1$, $0<\epsilon\le \min\{x,1-x\}$ y $B'=(y-\delta,y+\delta)$ para $0<y<1$, $0<\delta\le \min\{y,1-y\}$, $\max\{x-\epsilon,y-\delta\}<\min\{x+\epsilon,y+\delta\}$. En este caso tenemos
                    \[
                    (x-\epsilon,x+\epsilon)\cap(y-\delta,y+\delta)=\left(\frac{m_1+m_2}{2}-\epsilon_0,\frac{m_1+m_2}{2}+\epsilon_0\right)\in\mathcal{B}
                    \]
                    donde $m_1=\max\{x-\epsilon,y-\delta\}$, $m_2=\min\{x+\epsilon,y+\delta\}$ y $\epsilon_0 = \frac{m_2-m_1}{2}$.
        
                    \item Supongamos que $B=(x-\epsilon,x+\epsilon)$ para $0<x<1$, $0<\epsilon\le \min\{x,1-x\}$ y $B'=(1-\lambda,1]$ para $0<\lambda\le 1$, $1-\lambda<x+\epsilon$. Entonces tenemos
                    \[
                    (x-\epsilon,x+\epsilon)\cap(1-\lambda,1]=(1-\lambda,x+\epsilon)=(m_3-\epsilon_1,m_3+\epsilon_1)\in\mathcal{B}
                    \]
                    donde $m_3=\frac{x+\epsilon+1-\lambda}{2}$ y $\epsilon_1=\frac{x+\epsilon-1+\lambda}{2}$.
                \end{itemize}
        
                En los dos últimos casos también podemos argumentar que $B\cap B'$ es un abierto de $\mathcal{T}_u$. Así si $z\in B\cap B'$, teniendo en cuenta que \allowbreak $\{(z-\epsilon,z+\epsilon) \mid \epsilon>0\}$ es una base de entornos en $z$ de $\mathcal{T}_u$ podemos encontrar un $\epsilon>0$ suficientemente pequeño para $(z-\epsilon,z+\epsilon)\subset B\cap B'$.
        
                Finalmente, observemos que para los puntos en $(0,1]$, los elementos de $\mathcal{B}$ coinciden con una base de $\mathcal{T}_u|_{(0,1]}$. Por tanto, las bases de entornos en $x\in(0,1]$ para $\mathcal{T}_u|_{(0,1]}$ son bases de entornos en esos puntos para la topología $\mathcal{T}$. Se tienen entonces las siguientes bases de entornos para un punto $x\in X$ de la topología $\mathcal{T}$:
                \begin{align*}
                \mathcal{B}_0&=\{\{0\}\},\\
                \mathcal{B}_x&=\{(x-\epsilon,x+\epsilon) \mid 0<\epsilon\le \min\{x,1-x\}\} \quad \text{si } x\in(0,1),\\
                \mathcal{B}_1&=\{(1-\epsilon,1] \mid 0<\epsilon\le 1\},\\
                \mathcal{B}_x&=\{\{0,x\}\} \quad \text{si } x\in[2,3].
                \end{align*}
        
            \item Calcula la adherencia y el interior en $(X, \mathcal{T})$ de los siguientes subconjuntos $B_1 = \{0\}$ y $B_2 = X \cap \mathbb{Q}$.

                En primer lugar observamos que el conjunto $B_1$ es abierto ya que está en $\mathcal{B}$ y por tanto $\text{int}(B_1)=\{0\}$. Veamos ahora que $\text{cl}(B_1)=\{0\}\cup[2,3]$. Efectivamente, en primer lugar tenemos $0\in \text{cl}(B_1)$. Además $x\in[2,3]$ tenemos que el único elemento de la base de entornos en $x$ de $\mathcal{T}$ es $\{0,x\}$ y por tanto $B_1\cap\{0,x\}=\{0\}\ne\emptyset$. Luego $\{0\}\cup[2,3]\subset \text{cl}(B_1)$. Por otra parte $\{0\}\cup[2,3]=X\setminus(0,1]$ luego es un cerrado por ser $(0,1]$ un abierto. Como la clausura es el menor cerrado que contiene al conjunto deducimos que $\text{cl}(B_1)\subset\{0\}\cup[2,3]$. De la doble inclusión deducimos la igualdad $\text{cl}(B_1)=\{0\}\cup[2,3]$.

                Veamos ahora que $\text{int}(B_2)=\{0\}\cup([2,3]\cap\mathbb{Q})$. Observemos que para cada $x\in [2,3]\cap\mathbb{Q}$, $x\in\{0,x\}\subset\{0\}\cup([2,3]\cap\mathbb{Q})$, luego es entorno de todos sus puntos y por tanto $\{0\}\cup([2,3]\cap\mathbb{Q})$ es abierto. Como $\text{int}(B_2)$ es el abierto más grande contenido en $B_2$ tenemos $\{0\}\cup([2,3]\cap\mathbb{Q})\subset \text{int}(B_2)$. Por otra parte, si $x\in(0,1]\cap\mathbb{Q}$ es fácil comprobar que $x\notin \text{int}(B_2)$ ya que cualquier entorno básico de $x$ de la forma $(x-\epsilon,x+\epsilon)$ contiene números irracionales y por tanto nunca puede estar contenido en $B_2$. Concluimos por tanto que $\text{int}(B_2)=\{0\}\cup([2,3]\cap\mathbb{Q})$.
        
                Veamos ahora que $B_2$ es denso, es decir $\text{cl}(B_2)=X$. De hecho si consideramos $x\in(0,1)$, sea racional o irracional, tenemos que cualquiera de sus entornos básicos $(x-\epsilon,x+\epsilon)$, $0<\epsilon\le \min\{x,1-x\}$ contiene números racionales y por tanto $x\in \text{cl}(B_2)$. Análogamente si $x\in[2,3]$, sea racional o irracional, su entorno básico $\{0,x\}$ contiene al $0$ que es un número racional y por tanto $x\in \text{cl}(B_2)$. Concluimos entonces que $\text{cl}(B_2)=X$.
            
            \item ¿Es $(X, \mathcal{T})$ $T_1$? ¿Es $T_2$? ¿Es $1\text{AN}$? ¿Es $2\text{AN}$?

                Veamos que $(X,\mathcal{T})$ no es $T_1$. Efectivamente, consideremos los puntos $x=0$ y $y=2$. Cualquier entorno de $y=2$ contiene a $0$ y por tanto el espacio topológico no es $T_1$. Como no es $T_1$ no puede ser tampoco $T_2$.

                Veamos ahora que el espacio es 1AN. Para ello observemos que en el caso $x=0$ y $x\in[2,3]$ la base de entornos dada en el apartado a) es finita y por tanto numerable. Si $x\in(0,1)$ podemos considerar la base de entornos numerable:
                \[
                \mathcal{B}_x=\left\{\left(x-\frac{1}{n},x+\frac{1}{n}\right) \;\middle|\; n\in\mathbb{N}, 0<\frac{1}{n}\le \min\{x,1-x\}\right\}.
                \]
                Análogamente, para $x=1$ tenemos la base de entornos numerable:
                \[
                \mathcal{B}_1=\left\{\left(1-\frac{1}{n},1\right] \;\middle|\; n\in\mathbb{N}, 0<\frac{1}{n}\le 1\right\}.
                \]
                Por tanto concluimos que $(X,\mathcal{T})$ es 1AN.
        
                Veamos ahora que $(X,\mathcal{T})$ no es 2AN. Si $(X,\mathcal{T})$ fuese 2AN tendríamos que $([2,3],\mathcal{T}_{[2,3]})$ sería 2AN por ser esta una propiedad hereditaria. Pero $\mathcal{T}_{[2,3]}$ es la topología discreta en $[2,3]$ y esta topología sabemos que no es 2AN si el conjunto no es numerable.
            
            \item ¿Es $(X, \mathcal{T})$ conexo? ¿Es compacto?

                Es sencillo comprobar que $(X,\mathcal{T})$ no es conexo.

                \textbf{PRIMERA FORMA:} Observemos que el conjunto $Y=\{0\}\cup[2,3]$ es cerrado (visto en el apartado b), abierto (porque es entorno de todos sus puntos) e $Y\ne X$, $Y\ne\emptyset$.
        
                \textbf{SEGUNDA FORMA:} Denotemos por
                \[
                A=(0,1], \quad B=\bigcup_{x\in[2,3]}\{0,x\},
                \]
                entonces $A$ y $B$ son no vacíos, abiertos, disjuntos, $A\cup B=X$ luego $X$ no es conexo.
        
                Veamos ahora que $(X,\mathcal{T})$ no es compacto.
        
                \textbf{PRIMERA FORMA:} Si $(X,\mathcal{T})$ fuese compacto tendríamos que $([2,3],\mathcal{T}_{[2,3]})$ sería un compacto por ser $[2,3]\in\mathcal{C}_{\mathcal{T}}$ ($[2,3]=X\setminus[0,1]$ y $[0,1]\in\mathcal{T}$). Pero como ya hemos comentado en el apartado anterior $\mathcal{T}_{[2,3]}$ es la topología discreta en $[2,3]$ y esta topología sabemos que no es compacta si el conjunto es infinito.
        
                \textbf{SEGUNDA FORMA:} Consideramos el siguiente recubrimiento por abiertos (básicos) de $X$:
                \[
                X=(0,1]\cup\left(\bigcup_{x\in[2,3]}\{0,x\}\right)
                \]
                Si omitimos un abierto de esta familia, ya no es un recubrimiento de $X$, pues quedarían puntos sin cubrir. Así que no podemos extraer ningún subrecubrimiento finito y, por tanto, el espacio no es compacto.
        
            \item Se considera el conjunto $A = \{0\} \cup \left[2, 3\right]$. Determina si el espacio topológico cociente $(X/A, \mathcal{T}/A)$ es homeomorfo a $\left(\left[0, 1\right], \mathcal{T}_u|_{\left[0, 1\right]}\right)$, donde $\mathcal{T}_u$ es la topología usual en $\mathbb{R}$.

                Veamos que $(X/A,\mathcal{T}/A)$ no es homeomorfo a $([0,1],\mathcal{T}_u|_{[0,1]})$.
    
                \textbf{PRIMERA FORMA:} Veamos que $(X/A,\mathcal{T}/A)$ es un espacio topológico disconexo mientras que $([0,1],\mathcal{T}_u|_{[0,1]})$ es conexo por ser un intervalo en $\mathbb{R}$ y la conexión es un invariante topológico. De las afirmaciones anteriores solo resta comprobar que $(X/A,\mathcal{T}/A)$ no es conexo. Denotemos $p:X\rightarrow X/A$ la proyección canónica. Observemos que podemos escribir
                \[
                X/A=p(\{0\})\cup p((0,1])
                \]
                y tanto $p(\{0\})$ como $p((0,1])$ son abiertos, disjuntos y no vacíos ya que
                \[
                p^{-1}(p(\{0\}))=A\in\mathcal{T}, \quad p^{-1}(p((0,1]))=(0,1]\in\mathcal{T}, \quad p(\{0\})\cap p((0,1])=\emptyset
                \]
                También se puede razonar que $p(\{0\})$ y $p((0,1])$ son abiertos diciendo que se ha probado en teoría que si $A$ es abierto la aplicación $p:X\rightarrow X/A$ es abierta.
        
                \textbf{SEGUNDA FORMA:} También se puede argumentar que el espacio cociente $(X/A,\mathcal{T}/A)$ no es compacto, luego no puede ser homeomorfo a $([0,1],\mathcal{T}_u|_{[0,1]})$ que sí lo es. En efecto, un recubrimiento por abiertos de $X/A$ es
                \[
                X/A=p(\{0\})\cup\left(\bigcup_{n\in\mathbb{N}}p\left(\left(\frac{1}{n},1\right]\right)\right)
                \]
                y no se puede extraer ningún subrecubrimiento finito.
            
            \item Se considera el conjunto $E = \{0, 1\} \cup \left[2, 3\right]$. Determina si el espacio topológico cociente $(X/E, \mathcal{T}/E)$ es homeomorfo a $\left(\left]0, 1\right], \mathcal{T}_u|_{\left]0, 1\right]}\right)$, donde $\mathcal{T}_u$ es la topología usual en $\mathbb{R}$.

                Veamos que en este caso los espacios topológicos sí son homeomorfos. Consideremos la aplicación $f:(X,\mathcal{T})\rightarrow((0,1],\mathcal{T}_u|_{(0,1]})$ dada por
                \[
                f(x)=\begin{cases}
                x & \text{si } x\in(0,1],\\
                1 & \text{si } x\in\{0\}\cup[2,3].
                \end{cases}
                \]
                Claramente la aplicación $f$ es sobreyectiva.
        
                Observemos ahora que $A_1=(0,1]$ y $A_2=\{0\}\cup[2,3]$ son abiertos disjuntos. Además $f|_{A_1}=\text{Id}:(A_1,\mathcal{T}_{A_1})\rightarrow(A_1,\mathcal{T}_u|_{A_1})$ que es continua porque como ya hemos comentado en el apartado a) se tiene $\mathcal{T}_{A_1}=\mathcal{T}_u|_{A_1}$ y $f|_{A_2}=c_1:(A_2,\mathcal{T}_{A_2})\rightarrow(A_2,\mathcal{T}_u|_{A_2})$ es una aplicación constante y por tanto continua. El Lema del pegado nos garantiza entonces que la aplicación $f$ es continua.
        
                Veamos ahora que $f$ admite una inversa continua por la derecha y tendremos por un resultado de clase que $f$ es una identificación. Usando de nuevo que $\mathcal{T}_{A_1}=\mathcal{T}_u|_{A_1}$ tenemos que la inclusión $i:(A_1,\mathcal{T}_u|_{A_1})\rightarrow(X,\mathcal{T})$ es una aplicación continua. Además $f\circ i=\text{Id}_{A_1}$ de donde se tiene que $i$ es una inversa continua por la derecha de $f$ y por tanto $f$ es una identificación.
        
                Por otra parte es obvio que $R=R_f$ y así la aplicación inducida por $f$ en el cociente $\tilde{f}:(X/E,\mathcal{T}/E)\rightarrow((0,1],\mathcal{T}_u|_{(0,1]})$ es un homeomorfismo.
            
        \end{enumerate}
    \end{ejercicio}

    \begin{ejercicio}
        Un espacio topológico $(X, \mathcal{T})$ se dice $T_3$ si es $T_1$ y cumple que para cada subconjunto cerrado $C \in \mathcal{C}_{\mathcal{T}}$ y cada punto $x \in X \setminus C$ existen abiertos $A_1, A_2 \in \mathcal{T}$ tales que $C \subset A_1$, $x \in A_2$ y $A_1 \cap A_2 = \emptyset$.
        
            \begin{enumerate}[label=\alph*)]
                \item Sea $(X, \mathcal{T})$ un espacio topológico $T_1$. Demuestra que $(X, \mathcal{T})$ es $T_3$ si, y solo si, para cada punto $x \in X$ y cada abierto $U \in \mathcal{T}$ con $x \in U$, existe un abierto $V \in \mathcal{T}$ tal que $x \in V$ y $\overline{V} \subset U$.
    
                    \begin{itemize}[leftmargin=*]
                \item[$\implies$] Sea $x\in X$ y $U\in\mathcal{T}$ con $x\in U$. Comprobemos que existe un abierto $V\in\mathcal{T}$ tal que $x\in V$ y $\overline{V}\subset U$.
                
                Denotemos $C=X\setminus U$ que es un cerrado de $(X,\mathcal{T})$ tal que $x\in X\setminus C$. Por la definición de $T_3$ tenemos entonces que existen abiertos $A_1, A_2\in\mathcal{T}$ tales que $C\subset A_1$, $x\in A_2$ y $A_1\cap A_2=\emptyset$. Veamos que el abierto $V=A_2$ es el que buscamos. Efectivamente, en primer lugar tenemos que $x\in V$. Por otra parte tenemos que $V=A_2\subset X\setminus A_1$ que es un cerrado de $(X,\mathcal{T})$ y por tanto $\overline{V}\subset X\setminus A_1\subset X\setminus C=U$.
    
                \item[$\impliedby$] Sean $C\in\mathcal{C}_{\mathcal{T}}$ y $x\in X\setminus C$. Veamos que existen abiertos $A_1, A_2\in\mathcal{T}$ tales que $C\subset A_1$, $x\in A_2$ y $A_1\cap A_2=\emptyset$.
                
                Denotemos $U=X\setminus C\in\mathcal{T}$ que verifica $x\in U$. Por hipótesis sabemos que existe un abierto $V\in\mathcal{T}$ tal que $x\in V$ y $\overline{V}\subset U$. Veamos que $A_1=X\setminus\overline{V}$ y $A_2=V$ son dos abiertos que verifican la condición de $T_3$. Efectivamente observemos que en primer lugar $A_1, A_2\in\mathcal{T}$. Por otra parte $C=X\setminus U\subset X\setminus\overline{V}=A_1$, $x\in A_2$ y $A_1\cap A_2=(X\setminus\overline{V})\cap V\subset(X\setminus V)\cap V=\emptyset$.
            \end{itemize}
        
            \item Demuestra que $T_3$ es un invariante topológico.
            
                Tenemos que probar que si $f:(X,\mathcal{T})\rightarrow(Y,\mathcal{T}')$ es un homeomorfismo y $(X,\mathcal{T})$ es un espacio topológico $T_3$ entonces $(Y,\mathcal{T}')$ es también un espacio topológico $T_3$.
    
                Observemos en primer lugar que si $(X,\mathcal{T})$ es un espacio topológico $T_1$, por ser $T_1$ un invariante topológico, se tiene que $(Y,\mathcal{T}')$ es también un espacio topológico $T_1$.
        
                Consideremos ahora $C'$ un cerrado de $(Y,\mathcal{T}')$ e $y\in Y\setminus C'$. Denotemos $x=f^{-1}(y)$ y $C=f^{-1}(C')$. Observemos que $C\in\mathcal{C}_{\mathcal{T}}$ por ser $f$ continua y $x\in X\setminus C$. Por tanto como $(X,\mathcal{T})$ es un espacio topológico $T_3$ tenemos que existen abiertos $A_1, A_2\in\mathcal{T}$ tales que $C\subset A_1$, $x\in A_2$ y $A_1\cap A_2=\emptyset$. Tenemos entonces que $A_1'=f(A_1), A_2'=f(A_2)\in\mathcal{T}'$ ya que al ser $f$ un homeomorfismo es una aplicación abierta. Además como $C\subset A_1$ tenemos que $C'=f(C)\subset A_1'$, $y=f(x)\in A_2'$ y
                \[
                A_1'\cap A_2'=f(A_1)\cap f(A_2) \overset{f\text{ inyectiva}}{=} f(A_1\cap A_2) = f(\emptyset) = \emptyset.
                \]
        \end{enumerate}
    \end{ejercicio}
\end{document}